\scanpage{178}\mspage{33}%
% archive sheet number: 178
% Reader C: independent reading of scan 178, author page 33.
vérifier et d'expliciter comment les ``6 opérations'' habituelles sur les faisceaux (soit sur $X$, soit sur une réunion de strates de $X$) peuvent s'exprimer dans ce dictionnaire, dans le cas disons de strates non singulières (autrement, on aurait des difficultés avec les complexes dualisants, dont on aimerait qu'ils appartiennent à notre catégorie%
% A compact point is below and to the left of the faint horizontal dot row,
% approximately (949,661) in the whole image; its origin is uncertain.
\raisebox{-0.45ex}[0pt][0pt]{\parbox{0pt}{\figure{\begin{tikzpicture}[overlay]\fill (0.08,0) circle (0.35pt);\end{tikzpicture}}}}%
\unsure{.\,.\,.\,.\,.|\ldots}) \inserted{et de revérifier le formulaire connu de ces opérations}. Mais il apparaitra probablement que, pour bien pouvoir tout transcrire, il convient, plutôt \inserted{que} de regarder un diagramme du type
\diagram{\begin{tikzcd}[ampersand replacement=\&,column sep=1.4em,row sep=1.2em]
\cdot \arrow[r] \arrow[d] \& {}\raisebox{0pt}[0pt][0pt]{\parbox{0pt}{\figure{\begin{tikzpicture}[overlay]\fill (0.25,-0.15) ellipse (0.35pt and 0.25pt);\end{tikzpicture}}}} \\
{} \& {}
\end{tikzcd}}
% The zero-width point beside the empty upper-right diagram cell records
% the isolated mark near (1230,905), below-right of the arrowhead. Its
% origin is uncertain; it is not transcribed as a diagram object or label.
dans la catégorie homotopique formée avec la cat.\,. des ens. semi-simpliciaux, \inserted{de} regarder la catégorie des diagrammes d'ensembles $1/2$ simpliciaux, et à partir de là passer à une catégorie homotopique de fractions.

J'ai plus ou moins explicité dernièrement, en cogitant sur \deleted{\unsure{les|des|ces}}\unsure{des|les} fondements de ``topologie modérée'' (i.e. où on élimine d'emblée tous les phénomènes de sauvagerie), comment une stratification équisingulière\inserted{\deleted{(d'un espace \unsure{modéré} \unsure{compact})}}, à strates non singulières disons, d'un ``espace modéré'' \inserted{compact}, donne naissance canoniquement à un diagramme \deleted{\unsure{fon|f}} d'espaces,
% The cancelled parenthesis is a smaller interlinear addition above
% équisingulière, with a long enclosing/attachment stroke to its right end.
% compact is a separate interlinear insertion after the closing quotation
% marks of espace modéré and before the separately visible comma.
